\section*{NeurIPS Paper Checklist}

\begin{enumerate}

\item {\bf Claims}
    \item[] Question: Do the main claims made in the abstract and introduction accurately reflect the paper's contributions and scope?
    \item[] Answer: \answerYes{}
    \item[] Justification: The abstract and \S\ref{sec:introduction} advance four scoped claims, each backed by a specific theorem, table, or figure in the paper. The unifying message is that the \emph{simplest possible soft prompt}---training-free, freshly resampled per rollout, carrying no learned content---already reproduces the entropy-shaping signature and Pass@$N$ benefit of optimized methods, while admitting a clean theoretical account.

    (i) The per-layer RSP contribution decomposes exactly (Eq.~\eqref{eq:decomp}) into a real-token term and a contribution governed by a single attention-mass scalar $w_{r,i}^{(\ell)}$, whose fixed-gap upper bound is shown in Theorem~\ref{thm:decay} to shrink monotonically with KV cache growth, yielding an implicit explore-then-exploit dynamic without any scheduler (\S\ref{sec:converge}).

    (ii) Proposition~\ref{thm:minimax_directional_coverage} shows that isotropic Gaussians uniquely maximize worst-case directional variance at a fixed second-moment budget; combined with full support, this lets RSP reach top-$K$ entry regions inaccessible to any rank-preserving temperature transform. Empirically, RSP at $\tau=1.0$ produces Spearman $\rho=0.709$ against the $\tau=1.0$ baseline first-token distribution (vs $\rho=1$ for any temperature) while preserving $73.30\%$ Pass@1 (vs $1.42\%$ at the $\tau=3.0$ that would be needed to match RSP's distributional spread; Table~\ref{tab:rsp_vs_temp}).

    (iii) RSP shares the early-decoding entropy$\uparrow$/varentropy$\uparrow$ signature of LTPO, SoftCoT, and TTSV---three methods trained under \emph{different} objectives---supporting the interpretation that part of what trained \textit{soft prompts} deliver is a structural by-product of injection rather than of learned content (Figure~\ref{fig:entropy_first5}). Pass@$N$ accumulates only under \emph{independent} per-rollout resampling: at $N{=}16$, independent-seed RSP at $\tau{=}1.0$ reaches $92.80\%$ on MATH-500 and $36.67\%$ on AIME 2024, while a shared single-seed RSP shows no improvement over the temperature-only baseline (Figure~\ref{fig:pass_at_n}).

    (iv) The same input-side perturbation composes with DAPO's loss-side diversity preservation: DAPO\,+\,RSP attains a five-benchmark average of $54.36\%$ (vs DAPO's $53.20\%$ peak; $+1.16$\,pp at peak, $+3.10$\,pp at step~100), with the post-peak DAPO decline delayed (Table~\ref{tab:dapo_peak}, Figure~\ref{fig:dapo_curve}).

\item {\bf Limitations}
    \item[] Question: Does the paper discuss the limitations of the work performed by the authors?
    \item[] Answer: \answerYes{}
    \item[] Justification: Limitations are noted across the paper rather than in a single dedicated section, with future-work directions consolidated in \S\ref{sec:conclusion}. (a) The Pass@$N$ accumulation argument requires the task-side assumption $p_{\min}(x)>0$, explicitly flagged in \S\ref{sec:converge}; the isotropy and full-support arguments in \S\ref{sec:multipath} cover only the local branch-opening step. (b) The $+18$--$29$\,pp recovery on the format-mismatch rows of Table~\ref{tab:performance} is reported as an empirical observation, with the EOS-derailment hypothesis stated informally in \S\ref{sec:performance} rather than derived from any theorem in \S\ref{sec:why_rsp_works}. (c) Evaluation rigor: RSP length $L\in\{10,15,20\}$ and injection position are selected per model on the evaluation set without held-out validation, and Table~\ref{tab:performance} together with the DAPO\,+\,RSP run uses a single seed; multi-seed evaluation and held-out validation are listed as future work in \S\ref{sec:conclusion}~(iii)(iv) and discussed in item~\ref{checklist:significance}. (d) Formal-result scope: Theorem~\ref{thm:decay} guarantees only sequence-axis attenuation at fixed (layer, query, gap), with the layer-axis decay in Figure~\ref{fig:attention_mass} outside its envelope (\S\ref{sec:conclusion}~(ii)); Proposition~\ref{thm:minimax_directional_coverage} is a design criterion stated in input space and does not claim that isotropy survives transformer layers (\S\ref{sec:multipath} footnote). (e) The analysis is confined to RoPE-based decoders and mathematical reasoning; whether the same mechanism transfers to other position encodings (ALiBi, NoPE) and domains (code, open-ended reasoning) is unverified and listed as future work (\S\ref{sec:conclusion}~(i)).

\item {\bf Theory assumptions and proofs}
    \item[] Question: For each theoretical result, does the paper provide the full set of assumptions and a complete (and correct) proof?
    \item[] Answer: \answerYes{}
    \item[] Justification: The paper contains four formal theoretical results, each with explicit assumptions and a proof. (i) Equation~\eqref{eq:decomp} (head-output decomposition into a renormalized real-token term and an RSP-induced contribution) carries a one-line in-text derivation in \S\ref{sec:multipath}, using only the softmax-normalization identity under finite attention logits with $n\ge 1$ real and $L\ge 1$ RSP tokens. (ii) Theorem~\ref{thm:decay} (KV cache-growth attenuation) is proved inline in \S\ref{sec:converge}, assuming scaled dot-product attention with $n\ge 1$ real keys, $L\ge 1$ RSP keys, finite logits, and the pre-scale attention-logit gap $\Delta_i$ as defined in the statement. (iii) Proposition~\ref{thm:minimax_directional_coverage} (maximin directional coverage) is stated in \S\ref{sec:multipath} with its full proof in Appendix~\ref{app:why_random}, under the second-moment budget $\mathrm{tr}(\Sigma_D)\le\rho^2$ on zero-mean distributions over $\mathbb{R}^d$. (iv) The Pass@$N$ ensemble bound $\Pr(\mathrm{Pass@}N\text{ on }x)\ge 1-(1-p_{\min}(x))^N$ is stated in \S\ref{sec:converge} and proved formally in Appendix~\ref{app:passN_ensemble}, requiring (M1) the task-side condition $p_{\min}(x)>0$ and (M2) mutual independence of the $N$ rollouts. The supporting top-$K$ entry argument in \S\ref{sec:multipath}, which combines full support of the RSP law with a locally affine surrogate $\Delta_{ab}(\bar h)\approx\Delta_{ab}(0)+b_{ab}^\top\bar h$ around $\bar h=0$, is presented informally in the main text and written out formally in Appendix~\ref{app:topk_entry}.

    \item {\bf Experimental result reproducibility}
    \item[] Question: Does the paper fully disclose all the information needed to reproduce the main experimental results of the paper to the extent that it affects the main claims and/or conclusions of the paper (regardless of whether the code and data are provided or not)?
    \item[] Answer: \answerYes{}
    \item[] Justification: All hyperparameters required to reproduce the main experiments are disclosed in \S\ref{sec:setup}, \S\ref{sec:diversity}, \S\ref{sec:application}, and Appendices~\ref{app:setup}--\ref{app:dapo}. For the inference experiments, we report the decoding strategy (greedy for Tables~\ref{tab:performance}--\ref{tab:prior_comparison}; $\tau=1.0$ sampling for the diversity and Pass@$N$ experiments in \S\ref{sec:diversity}), the RSP length $L\in\{10,15,20\}$ together with its per-model selection protocol on the evaluation set (Tables~\ref{tab:performance}--\ref{tab:prior_comparison}) and the fixed $L=10$, \textit{suffix} configuration used in the mechanism analyses (\S\ref{sec:entropy}--\S\ref{sec:diversity}), the injection position grid \{\textit{prefix}, \textit{suffix}\} (with per-position breakdown in Appendix~\ref{app:ablation}, and the exact prompt templates for all three positions \{\textit{prefix}, \textit{infix}, \textit{suffix}\} across the ChatML, LLaMA, and raw-text formats in Appendix~\ref{app:injection_positions}), the number of samples per problem (16 for the token-level and Pass@$N$ analyses), the number of RSP seeds (16 for the token metrics in Table~\ref{tab:rsp_vs_temp}), and the trajectory-level setting (64 problems with 300 independent trajectories per condition encoded by BGE-M3 \citep{chen2024bge} and clustered by DBSCAN \citep{ester1996density}). The unified answer-extraction pipeline (Appendix~\ref{app:extraction}) and per-benchmark generation-length budgets are also reported. For the DAPO + RSP training experiments, Appendix~\ref{app:dapo} discloses the optimizer, the prompt batch ($1{,}024$ prompts), the per-prompt rollout count ($G=8$), the rollout temperature ($\tau=1.0$), the four PPO mini-batch updates of $256$ prompts each, the total number of training steps ($100$), the asymmetric clipping range $(\varepsilon_{\text{low}}, \varepsilon_{\text{high}})=(0.2, 0.28)$, the dual-clipping safeguard, the suffix RSP length used during rollouts ($L=20$), the rollout-only application protocol (no RSP at evaluation), and the dataset (a Level-3--5 subset of MATH with $\sim$8.5K prompts). The five-benchmark evaluation protocol (GSM8K, MATH-500, College Math, Minerva Math, AIME 2024) specifies Avg@32 at $\tau=1.0$ for AIME 2024 and greedy decoding elsewhere. The inference batch sizes are also disclosed (16 for LLaMA-3.1-8B-Instruct, 32 for the Qwen2.5-Math 1.5B variants).

\item {\bf Open access to data and code}
    \item[] Question: Does the paper provide open access to the data and code, with sufficient instructions to faithfully reproduce the main experimental results, as described in supplemental material?
    \item[] Answer: \answerYes{}
    \item[] Justification: All evaluation data and base models used in the paper are publicly available and cited at first use. Beyond the description in the main text (Eq.~\eqref{eq:rsp_def}, \S\ref{sec:setup}, Appendix~\ref{app:dapo}), an anonymized implementation of the RSP injection harness and the DAPO\,+\,RSP training pipeline is included as supplementary material, with a README detailing the run commands for the inference experiments in \S\ref{sec:performance}--\S\ref{sec:diversity} and the DAPO\,+\,RSP training in \S\ref{sec:application}. Pretrained DAPO and DAPO\,+\,RSP checkpoints across training steps 10--100 will be released at the camera-ready stage to allow reviewers and downstream users to reproduce the main DAPO\,+\,RSP results in \S\ref{sec:application} without re-running the training stage.

\item {\bf Experimental setting/details}
    \item[] Question: Does the paper specify all the training and test details (e.g., data splits, hyperparameters, how they were chosen, type of optimizer) necessary to understand the results?
    \item[] Answer: \answerYes{}
    \item[] Justification: Beyond the hyperparameter values disclosed for reproducibility (item~4), this item details the data splits, hyperparameter selection protocol, optimizer choices, and prior-work hyperparameters. \textbf{Data splits.} The inference experiments (\S\ref{sec:performance}--\S\ref{sec:diversity}) involve no training and use the standard public test sets of MATH-500, GSM8K, and AIME 2024. The DAPO + RSP training (\S\ref{sec:application}) uses a Level-3--5 subset of MATH (\texttt{simplelr\_math\_35}, $\sim$8.5K prompts) for training and evaluates on five held-out benchmarks (GSM8K, MATH-500, College Math, Minerva Math, AIME 2024) under the per-benchmark sampling protocol of Appendix~\ref{app:dapo} (Table~\ref{tab:dapo_eval_protocol}). \textbf{Selection protocol.} The RSP token count $L\in\{10,15,20\}$ is selected per (model, benchmark) on the evaluation set, with the per-row choice in Appendix~\ref{app:setup} (Table~\ref{tab:token_count}); the injection position is reported as a grid (\textit{prefix}, \textit{infix}, \textit{suffix}) in Appendix~\ref{app:setup} (Table~\ref{tab:performance_full}) rather than tuned, and the mechanism analyses in \S\ref{sec:entropy}--\S\ref{sec:diversity} fix $L=10$ and \textit{suffix} to insulate them from this selection effect. The absence of held-out validation is acknowledged in item~2 and \S\ref{sec:conclusion}~(iv). \textbf{Optimizer.} DAPO + RSP training uses AdamW with a constant learning rate of $5\times10^{-7}$ and no warmup; total training is approximately $10$ epochs with checkpoints saved every $10$ steps, as listed in Appendix~\ref{app:dapo} (Table~\ref{tab:dapo_hp}). \textbf{Prior work hyperparameters.} The reproduced TTSV, SoftCoT, and LTPO baselines used in Table~\ref{tab:prior_comparison} are documented with their full training and evaluation hyperparameters in Appendix~\ref{app:prior_hp}: TTSV (prefix length, AdamW, batch size, gradient accumulation, learning-rate schedule, per-model learning rates), LTPO (per-benchmark token count, optimization steps, top-$k$, sigma and decay, learning rate; Table~\ref{tab:hp_ltpo}), and SoftCoT (projection-module training schedule, thought-token counts during training and evaluation, batch size, gradient accumulation, and per-model learning rates).

\item {\bf Experiment statistical significance}\label{checklist:significance}
    \item[] Question: Does the paper report error bars suitably and correctly defined or other appropriate information about the statistical significance of the experiments?
    \item[] Answer: \answerNo{}
    \item[] Justification: Variability is reported for the experiments most directly affected by sampling stochasticity. Table~\ref{tab:rsp_vs_temp} reports Pass@1 as mean $\pm$ standard deviation over 16 samples per problem, with the token-level metrics additionally averaged over 16 independent RSP seeds. Figure~\ref{fig:entropy_first5} draws shading at $\pm 1$ SEM across MATH-500 problems. The Pass@$N$ analysis in Figure~\ref{fig:pass_at_n} uses $N=16$ samples per problem on MATH-500 and AIME 2024. The trajectory-level analysis (Table~\ref{tab:semantic_metrics}) is computed over 64 problems with 300 independent trajectories per condition and additionally reports the per-problem majority (RSP yields larger pairwise distance than baseline on 56 of 64 problems). On the DAPO + RSP evaluation side, AIME 2024 uses Avg@32 at $\tau=1.0$ to reduce competition-set variance. By contrast, Table~\ref{tab:performance} and Table~\ref{tab:prior_comparison} report single-run greedy-decoded accuracies without bootstrap confidence intervals or paired tests, and the DAPO + RSP training (Table~\ref{tab:dapo_peak}, Figure~\ref{fig:dapo_curve}) is run with a single training seed. Multi-seed evaluation across both inference and training, particularly on the small competition sets (AIME 2024, AMC 2023) where the per-problem variance is highest, is identified as future work in \S\ref{sec:conclusion}~(iii).

\item {\bf Experiments compute resources}
    \item[] Question: For each experiment, does the paper provide sufficient information on the computer resources (type of compute workers, memory, time of execution) needed to reproduce the experiments?
    \item[] Answer: \answerYes{}
    \item[] Justification: All inference experiments (Tables~\ref{tab:performance}--\ref{tab:semantic_metrics}, Figures~\ref{fig:entropy_first5}--\ref{fig:pass_at_n}) run on a single NVIDIA A6000 40GB GPU with greedy decoding (Appendix~\ref{app:setup}). DAPO + RSP training (\S\ref{sec:application}) runs on $4\times$ NVIDIA B200 GPUs for approximately $12$ hours per run (Appendix~\ref{app:dapo}, Table~\ref{tab:dapo_hp}, hardware and wall-clock rows).

\item {\bf Code of ethics}
    \item[] Question: Does the research conducted in the paper conform, in every respect, with the NeurIPS Code of Ethics \url{https://neurips.cc/public/EthicsGuidelines}?
    \item[] Answer: \answerYes{}
    \item[] Justification: The work uses publicly available math reasoning benchmarks and pretrained language models, involves no human subjects or personal data, and does not introduce a release that bypasses standard model-safety considerations. We have reviewed the NeurIPS Code of Ethics and identify no conflicts.

\item {\bf Broader impacts}
    \item[] Question: Does the paper discuss both potential positive societal impacts and negative societal impacts of the work performed?
    \item[] Answer: \answerYes{}
    \item[] Justification: Appendix~\ref{app:broader_impacts} discusses both directions: positive impact via more sample-efficient RL post-training of reasoning LLMs, and a potential downside that broadening the effective output support also increases reachability of undesirable trajectories, which we argue should keep RSP coupled with existing safety filtering rather than replacing it.

\item {\bf Safeguards}
    \item[] Question: Does the paper describe safeguards that have been put in place for responsible release of data or models that have a high risk for misuse (e.g., pre-trained language models, image generators, or scraped datasets)?
    \item[] Answer: \answerNA{}
    \item[] Justification: The DAPO and DAPO\,+\,RSP checkpoints (to be released at the camera-ready stage, item~5) are math-domain fine-tunes of the publicly available Qwen2.5-Math-7B \citep{yang2024qwen25math} on a Level-3--5 subset of MATH and do not extend the base model's capabilities beyond mathematical reasoning, so they inherit the upstream safety profile without introducing new misuse vectors. RSP itself is a training-free generation-time perturbation applied to existing pretrained models and adds no separate artifact requiring safeguards. The camera-ready checkpoint release will be accompanied by a model card stating the intended use (mathematical reasoning research) and the upstream license under which the base model is distributed.

\item {\bf Licenses for existing assets}
    \item[] Question: Are the creators or original owners of assets (e.g., code, data, models), used in the paper, properly credited and are the license and terms of use explicitly mentioned and properly respected?
    \item[] Answer: \answerYes{}
    \item[] Justification: All datasets (MATH-500 \citep{lightman2024lets}, GSM8K \citep{cobbe2021training}, AIME 2024, AMC 2023, OlympiadBench \citep{he2024olympiad}, Minerva Math \citep{lewkowycz2022minerva}, GaoKao 2023-EN, College Math) and pretrained models (Qwen2.5-Math \citep{yang2024qwen25math}, LLaMA-3.1 \citep{grattafiori2024llama}) are publicly available and cited at first use. The training pipeline builds on SimpleRL-Zoo \citep{zeng2025simplerl} and VeRL \citep{sheng2024hybridflow}, also cited. Each asset is used under its original publicly stated terms.

\item {\bf New assets}
    \item[] Question: Are new assets introduced in the paper well documented and is the documentation provided alongside the assets?
    \item[] Answer: \answerYes{}
    \item[] Justification: Two new assets are introduced to support reproducibility (item~5). (i) An anonymized implementation of the RSP injection harness and the DAPO\,+\,RSP training pipeline is included as supplementary material, accompanied by a README that documents the environment, run commands for the inference experiments in \S\ref{sec:performance}--\S\ref{sec:diversity}, and the DAPO\,+\,RSP training entry point used in \S\ref{sec:application}. (ii) Pretrained DAPO and DAPO\,+\,RSP checkpoints across training steps 10--100 will be released at the camera-ready stage, accompanied by a model card stating the intended use (mathematical reasoning research), a reference to the training data used (the Level-3--5 subset of MATH from SimpleRL-Zoo's publicly available \texttt{simplelr\_math\_35}, $\sim$8.5K prompts; the dataset itself is not redistributed), and the upstream Qwen2.5-Math-7B \citep{yang2024qwen25math} license that both checkpoints inherit. No new dataset is introduced and no human subjects are involved, so no consent procedure applies. Limitations of the released artifacts (single-seed training, RoPE/math-only scope) are documented in item~2 and \S\ref{sec:conclusion}.

\item {\bf Crowdsourcing and research with human subjects}
    \item[] Question: For crowdsourcing experiments and research with human subjects, does the paper include the full text of instructions given to participants and screenshots, if applicable, as well as details about compensation (if any)?
    \item[] Answer: \answerNA{}
    \item[] Justification: The paper involves no crowdsourcing and no research with human subjects. All evaluations use static, publicly available math reasoning benchmarks.

\item {\bf Institutional review board (IRB) approvals or equivalent for research with human subjects}
    \item[] Question: Does the paper describe potential risks incurred by study participants, whether such risks were disclosed to the subjects, and whether Institutional Review Board (IRB) approvals (or an equivalent approval/review based on the requirements of your country or institution) were obtained?
    \item[] Answer: \answerNA{}
    \item[] Justification: As noted in item~14, the paper involves no human subjects, so IRB review is not applicable.

\item {\bf Declaration of LLM usage}
    \item[] Question: Does the paper describe the usage of LLMs if it is an important, original, or non-standard component of the core methods in this research? Note that if the LLM is used only for writing, editing, or formatting purposes and does \emph{not} impact the core methodology, scientific rigor, or originality of the research, declaration is not required.
    \item[] Answer: \answerNA{}
    \item[] Justification: Pretrained LLMs (Qwen2.5-Math, LLaMA-3.1) appear in the paper as the experimental subject of analysis, not as a non-standard component of the methods themselves. The methods (RSP definition, attention decomposition, DAPO + RSP training) do not rely on auxiliary LLMs in any non-standard or originality-affecting role.

\end{enumerate}
